\documentclass[12pt]{report}
\usepackage[utf8]{inputenc}
\usepackage{tcolorbox}
\usepackage{tikz-cd}
\usepackage{multirow}
\usepackage{epsf}
\usepackage{tikz}
\usepackage{cancel}
\usepackage{hyperref}
\usepackage{marvosym}
\usepackage{tabularx}
\usepackage{booktabs}
\usepackage{float}
\usepackage{graphicx} % Required for inserting images
\usepackage{amsmath}
\usepackage{mathtools, amssymb}
\usepackage{enumitem}
\usepackage{amsthm}
\usepackage{amssymb}
\usepackage{xcolor}
\usepackage{array}
\usepackage{comment}
\usepackage{textgreek}
%\usepackage{evan}
\theoremstyle{definition}


\newtheorem*{theorem*}{Theorem}

%% this allows for theorems which are not automatically numbered

\newtheorem{thm}{Theorem}[section]
\newtheorem{lem}[thm]{Lemma}
\newtheorem{prop}[thm]{Proposition}
\newtheorem{claim}[thm]{Claim}
\newtheorem{remark}[thm]{Remark}
\newtheorem{example}[thm]{Example}
\newtheorem{fact}{Fact}
\newtheorem{corollary}[thm]{Corollary}
\newtheorem{defn}[thm]{Definition}
\newtheorem*{exercise}{Exercise}
\newtheorem*{motivation}{Motivation}
\newtheorem*{observation}{Observation}
\newtheorem*{problem}{Problem}
\newtheorem*{conjecture}{Conjecture}
\newtheorem*{notation}{Notation}
\newtheorem*{digression}{Digression}


%% Bunch of macros

%General
\newcommand{\bb}[1]{\mathbb{#1}}
\renewcommand{\cal}[1]{\mathcal{#1}}
\newcommand{\cyc}[1]{\mathbb{Z}/#1\mathbb{Z}}
\renewcommand{\vec}[1]{\mathbf{#1}}
\newcommand{\diff}[2]{\frac{\mathrm{d}#1}{\mathrm{d}#2}}
\newcommand{\pdiff}[2]{\frac{\partial#1}{\partial#2}}
\newcommand{\Int}[4]{\int\limits_{#1}^{#2}#3 \expandafter\ifx\expandafter\relax
  \detokenize{#4}\relax\else\textnormal{d#4}\fi} %if the argument for variable that we're integrating wrt is empty, this makes it so that there's no `d'
\renewcommand{\l}[1]{\left#1}
\renewcommand{\r}{\right}
\DeclareMathOperator{\img}{Image}
\DeclareMathOperator{\Frac}{Frac}
\newcommand{\pre}[1]{#1^{\textrm{pre}}}
\newcommand{\ovl}[1]{\overline{#1}}

%Algebra & Category Theory
\DeclareMathOperator{\id}{id}
\DeclareMathOperator{\chara}{char}
\DeclareMathOperator{\aut}{Aut}
\DeclareMathOperator{\inn}{Inn}
\DeclareMathOperator{\out}{Out}
\DeclareMathOperator{\stab}{Stab}
\newcommand{\gl}[2]{\text{GL}(#1, \mathbb{#2})}
\DeclareMathOperator{\Span}{Span}
\DeclareMathOperator{\End}{End}
\DeclareMathOperator{\ann}{Ann}
\DeclareMathOperator{\cl}{Cl}

%Analysis
\DeclareMathOperator{\im}{Im}
\DeclareMathOperator{\re}{Re}
\newcommand{\Oint}[4]{\oint\limits_{#1}^{#2}#3 \expandafter\ifx\expandafter\relax
  \detokenize{#4}\relax\else\textnormal{d#4}\fi} %if the argument for variable that we're integrating wrt is empty, this makes it so that there's no `d'
\newcommand{\norm}[1]{\left\lVert#1\right\rVert}


%Geometry/Topology
%\renewcommand{\norm}[1]{\left\lVert#1\right\rVert}
\DeclareMathOperator{\In}{Int}
\DeclareMathOperator{\conv}{conv}



%document format
\setlength{\parindent}{0pt}
\setlength{\parskip}{8pt}
\usepackage[left=2cm, right=2cm, top=1cm, bottom=2cm]{geometry}

\title{120KT Lecture Notes}
\author{moksh}
\begin{document}
\maketitle
\chapter{notes}
\section{Class 1}
\begin{defn}[Manifold]
A manifold of dimension $n$, or $n$-manifold is a 2nd countable Hausdorff topological space, $\cal{M}$, for which there exists a family $\{(V_{\lambda}, \phi_{\lambda})\}_{\lambda \in \lambda}$ such that $\{U_{\lambda}\}_{\lambda \in \Lambda} \subseteq \cal{T}_{\cal{M}}$, $\bigcup_{\lambda \in \Lambda} U_{\lambda} = \cal{M}$, and that each $\phi_{\lambda}$ is a homeomorphism between $U_{\lambda}$ and some open subset of $\bb{R}^n$.
\end{defn}

\begin{defn}[Coordinate chart]
Each $(U_{\lambda}, \phi_{\lambda})$ is called a coordinate chart of $\cal{M}$. 
The set of all coordinate charts is called an \emph{atlas} of $\cal{M}$.
\end{defn}

\begin{defn}[Transition map]
Let $(U_\alpha, \phi_\alpha), (U_\beta, \phi_\beta), \{\alpha, \beta\}\subseteq \Lambda$ be coordinate charts with $U_{\alpha} \cap U_{\beta} \neq \emptyset$.
 
Then we define the transition map $\tau_{\alpha, \beta}: \phi_{\alpha}(U_{\alpha} \cap U_{\beta}) \to \phi_{\beta}(U_{\alpha} \cap U_{\beta})$ as $\phi_{\beta} \circ \phi_{\alpha}^{-1}$.
\end{defn}
Clearly, $\tau_{\alpha, \beta}$ is a homeomorphism.

In the context of manifolds, equivalence is considered up to homeomorphism.


\begin{example}[Examples of manifolds]
Some examples:

\begin{enumerate}
\item $\bb{R}^n$ is an $n$-manifold.
\item Any open subset of $\bb{R}^n$ is an $n$-manifold.
\item $S^n = \{x \in \bb{R}^{n+1} : \norm{x} = 1\}$ is an $n$-manifold 
\item $\bb{RP}^n$ is an $n$-manifold, where $\bb{RP}^n$ is the quotient space obtained by identifying antipodal points on $\bb{S}^n$.
\end{enumerate}
\end{example}

\begin{exercise}
Prove that $S^n$ is an $n$-manifold.
\end{exercise}
\begin{proof}
\ref{snmanifold}
\end{proof}

\begin{exercise}
Prove that $\bb{RP}^n$ is an $n$-manifold.
\end{exercise}
\begin{proof}
\ref{rpnmanifold}
\end{proof}
\begin{defn}[Product Manifold]
Let $(\cal{M}, \{U_\lambda, \phi_\lambda\}_{\Lambda})$ and $(\cal{M}, \{U_\gamma, \phi_\gamma\}_{\Gamma})$ be $m$ and $n$ dimensional manifolds. 
Then, their product structure defines a manifold.
\end{defn}

\begin{exercise}
Prove this.

\begin{proof}
\ref{productmanifold}
\end{proof}
\end{exercise}

Then, we define the $n$-dimensional torus, denoted by $T^n$, as $\underbrace{S^1 \times S^1 \times  \dots \times S^1}_{n \text{ times}} = T^n$.

\begin{defn}[Sub-manifold]
Let $\cal{M}$ be an $n$-manifold, and let $\cal{K} \subseteq \cal{M}$, and let $k \leq n$.

We say $\cal{K}$ is a $k$-dimensional submanifold of $\cal{M}$, if $\cal{K}$ equipped with atlas $\{(U_{\lambda} \cap \cal{K}, \phi_{\lambda}\restriction_{K})\}_{\lambda \in \Lambda}$ is a $k$-manifold.
\end{defn}

If $\cal{K}$ is a submanifold of a manifold $\cal{M}$, we define its codimension as $n-k$. 
\begin{example}
$S^n$ is a codimension $1$ submanifold of $\bb{R}^{n+1}$.
\end{example}

\begin{defn}[Manifolds with boundary]
The usual definition of a manifold, except that some neighbourhoods are permitted to be homeomorphic to the upper half space, $\bb{H}^n$, which is the subset of $\bb{R}^n$ with first coordinate nonnegative. 
\end{defn}


\begin{defn}[Boundary \& Interior]
The boundary of a manifold $\cal{M}$ is denoted by $\partial \cal{M}$, and is defined as $\partial \cal{M} = \{x \in \cal{M}: x \in U_{\alpha} \Rightarrow U_{\alpha} \cong  \bb{H}^n\}$

Similarly, the interior of a manifold $\cal{M}$ is denoted by $\In \cal{M}$, and is defined as $\partial \cal{M} = \{x \in \cal{M}: x \in U_{\alpha} \Rightarrow U_{\alpha} \cong  U \in \cal{T}_{\bb{R}^n}\}$
\end{defn}
\begin{prop}
This notion is well defined, i.e. if for some point $x \in \cal{M}$, some open patch containing $x$ is homeomorphic to an open subset of $\bb{R}^n$, then any open patch containing $x$ is homeomorphic to an open subset of $\bb{R}^n$. 

Similarly, if for some point $x \in \cal{M}$, some open patch containing $x$ is homeomorphic to $\bb{H}^n$ , then any open patch containing $x$ is homeomorphic to an open subset of $\bb{R}^n$. 

\end{prop}
\begin{proof}
\ref{manifoldwboundary}
\end{proof}

%Mention the result she mentioned that boundary, interior are invariant regardless of coordinate chart chosen.
%Closed manifold Compact manifold, Open manifold, Simply Connected, Path connected, etc etc

\begin{defn}[Orientability]
A manifold is orientable iff $H_n(\cal{M}, \bb{Z}) \cong \bb{Z}$. 

\end{defn}


A manifold is orientable iff there can be defined a continuous normal vector field on it, when it is embedded in $\bb{R}^m$.

An orientation is just a choice of generator for its highest homology group.

$2$-manifolds are called surfaces.

$\Sigma_{g, b}$ are orientable surfaces of genus $g$, and \# int $D^2$ removed is $b$. 

$\Sigma_{0, 0}: S^2$.

$\Sigma_{0, 1}: D^2$.

$\Sigma_{0, 2}$: Annulus

$\Sigma_{0, 3} S^2$: Pair of pants

$\Sigma_{1, 0} T^2$.

Such manifolds are always compact, and have boundary for $g>0$.

\begin{defn}[Embedding]
Let $\cal{M}$, $\cal{N}$ be manifolds with $\dim(\cal{M}) \leq \dim(\cal{M})$. We say that $\cal{M}$ is embedded in $\cal{N}$ iff it is homeomorphic to some sub-manifold of $\cal{N}$, by some map $f$.
Clearly $f$ is injective, and we denote it by $f: \cal{M} \hookrightarrow \cal{N}$.
\end{defn}

Knot theory is the study of embeddings of $S^p$ in $S^n$ or $\bb{R}^n$, with $p< n$.

\chapter{Exercises}

\begin{prop}[$S^n$ is a manifold]
\label{snmanifold}
\end{prop}
\begin{proof}
By the stereographic projection, $\bb{R^n} \cong S^n \setminus (0, 0, \dots, 1)$, and by composing the stereographic projection with reflection in the $x_n$ plane, we get $\bb{R^n} \cong S^n \setminus (0, 0, \dots, -1)$.

Hence, these two sets are an explicit example of a collection of surface patches for $S^n$, which proves it is a manifold.
\end{proof}

\begin{prop}[$\bb{RP}^n$ is a manifold]
\label{rpnmanifold}
\end{prop}
\begin{proof}
\textcolor{red}{To be done.}
\end{proof}

\begin{prop}[Product manifolds are well defined]
\label{productmanifold}
\end{prop}
\begin{proof}
\textcolor{red}{To be done.}
\end{proof}

\begin{prop}[The boundary of a manifold is well defined]
\label{manifoldwboundary}
\end{prop}
\begin{proof}
\textcolor{red}{To be done.}
\end{proof}
\end{document}